%% Suspension arrière d'une moto électrique, vue de profil (schéma de principe).
%% Châssis 1 (vert, hachuré = référence), bras oscillant 2 (bleu) articulé en D sur le châssis,
%% combiné ressort-amortisseur : tige 3 (orange) articulée en B sur le bras, corps 4 (rouge)
%% articulé en C sur le châssis, la tige coulissant dans le corps suivant (BC).
%% Les deux flèches rouges sont les actions du ressort COMPRIMÉ (norme non chiffrée : à calculer).
\documentclass[tikz,border=4pt]{standalone}
\usepackage{liaisons}
\usetikzlibrary{backgrounds}
\begin{document}
\begin{tikzpicture}[x=1cm,y=1cm,
  piece/.style={line width=1.6pt, line cap=round, line join=round},
  fin/.style={line width=0.7pt},
  vect/.style={-{Stealth[length=6pt]}, line width=1.2pt, solideA},
  cote/.style={line width=0.6pt, solideD, {Stealth[length=4pt]}-{Stealth[length=4pt]}},
  traj/.style={line width=0.6pt, black!55, dash pattern=on 3pt off 2pt}]

\useasboundingbox (-0.4,-1.0) rectangle (8.1,4.5);
\coordinate (A) at (6.6,1.1);      % axe de la roue arrière
\coordinate (B) at (4.4,1.9);      % tige 3 / bras 2
\coordinate (C) at (5.35,3.65);    % corps 4 / châssis 1
\coordinate (D) at (2.6,1.5);      % bras 2 / châssis 1
\def\abc{61.5}                     % inclinaison de (BC) sur l'horizontale (≈ 60°)
\def\lbc{1.99}                     % longueur BC dans le dessin

%% ---------------- châssis 1 (référence, hachuré) --------------------------
\draw[piece, solideE] (D) -- (2.75,2.85) -- (3.95,3.35) -- (C) -- (4.95,2.85) -- (3.35,2.35) -- cycle;
\foreach \i in {0,1,...,6} {\pgfmathsetmacro\xx{3.95+1.40*\i/6}\pgfmathsetmacro\yy{3.35+0.30*\i/6}%
  \draw[line width=0.45pt, solideE] (\xx,\yy) -- (\xx-0.10,\yy+0.20);}
\node[font=\small, text=solideE, anchor=east] at (2.65,3.05) {\textbf{1} châssis};

%% ---------------- roue arrière ---------------------------------------------
\draw[line width=0.8pt, black!40] (A) circle (0.8);
\draw[line width=0.8pt, black!40] (A) circle (0.25);

%% ---------------- bras oscillant 2 -----------------------------------------
\draw[fin, solideB] (D) -- (B) -- (A);
\draw[piece, solideB] (D) -- (A);
\node[font=\small, text=solideB, sloped, anchor=south, inner sep=3pt] at (4.85,1.32)
     {\textbf{2} bras oscillant};

%% ---------------- trajectoire de A dans le mouvement 2/1 -------------------
\draw[traj, {Stealth[length=5pt]}-{Stealth[length=5pt]}]
  ([shift=(-20:4.02)]D) arc (-20:8:4.02);
\node[font=\scriptsize, text=black!55, anchor=east] at (7.7,-0.18) {trajectoire de $A$ (2/1)};

%% ---------------- combiné ressort-amortisseur ------------------------------
\begin{scope}[shift={(B)}, rotate=\abc]
  \draw[piece, solideD] (0,0) -- (0.95,0);                                   % tige 3
  \draw[line width=1.2pt, solideA, fill=white] (0.72,-0.17) rectangle (\lbc,0.17);  % corps 4
  \draw[line width=0.7pt, solideD] (0.52,-0.11) -- (0.90,-0.11)
                                   (0.52,0.11)  -- (0.90,0.11);              % glissière 3/4
  \draw[cote] (0,-0.85) -- (\lbc,-0.85);
\end{scope}
\draw[line width=0.5pt, solideD] (5.72,2.40) -- (6.30,2.45);
\node[font=\small, text=solideD, anchor=west, inner sep=2pt] at (6.30,2.45) {$BC=310$ mm};
\node[font=\small, text=solideD, anchor=east] at (4.25,2.45) {\textbf{3} tige};
\node[font=\small, text=solideA, anchor=west] at (6.25,3.45) {\textbf{4} corps};
\node[font=\small\itshape, text=black!60, anchor=north west] at (-0.35,0.95)
     {ressort non représenté};

%% ---------------- angle de 60° sur l'horizontale ---------------------------
\draw[line width=0.5pt, dash pattern=on 2.5pt off 2pt, black!60] (B) -- ++(1.0,0);
\draw[line width=0.6pt, solideD] ([shift=(0:0.62)]B) arc (0:\abc:0.62);
\node[font=\small, text=solideD] at ([shift=(24:0.72)]B) {$60^{\circ}$};

%% ---------------- actions du ressort comprimé -------------------------------
\draw[vect] (C) -- ++(\abc:0.70)
     node[anchor=west, inner sep=2pt, font=\scriptsize, text=solideA]
     {$\vec{F}_{\text{ressort}\rightarrow 4}$};
\draw[vect] (B) -- ++(\abc+180:0.70)
     node[anchor=north, inner sep=2pt, font=\scriptsize, text=solideA]
     {$\vec{F}_{\text{ressort}\rightarrow 3}$};

%% ---------------- liaisons pivot --------------------------------------------
\pic[s1=solideB, s2=solideE, liens=false] at (D) {pivot bout};
\pic[s1=solideD, s2=solideB, liens=false] at (B) {pivot bout};
\pic[s1=solideA, s2=solideE, liens=false] at (C) {pivot bout};

%% ---------------- points ------------------------------------------------------
\foreach \p/\pos in {A/{south east}, B/{south west}, C/{north west}, D/{north east}} {
  \fill (\p) circle (1.5pt);
  \node[font=\small\bfseries, anchor=\pos, inner sep=3pt] at (\p) {$\p$};}
\repere{(0.15,-0.7)}{x}{y}{1}
\end{tikzpicture}
\end{document}
